\section{MVDR 接收波束的全局最优性证明}

\noindent\textbf{定理。}考虑论文中的问题 $P0$。设 $P0$ 存在全局最优解，且感知信道满足式~\eqref{eq:A}，$\sigma^2>0$，$p_j(t)\ge 0$，$\bm a_r(\theta_i(t))\ne\bm 0$，$\omega_2\ge 0$。感知发射波束可以按 $[0,1]$ 内的系数缩放，且不存在未列出的发射功率下限、波束硬件约束或其他依赖波束的约束。感知发射、卸载发射和接收波束的定义域均为复向量空间 $\mathbb C^{N\times1}$。
则 $P0$ 至少存在一个全局最优解，其每个 UAV 在每个时隙的接收波束均为式~\eqref{eq:mvdr} 所给的单位范数 MVDR 接收波束。

\medskip
\noindent\textbf{证明。}任取问题 $P0$ 的可行解 $\tilde{\bm\zeta}$，保留其中所有时隙的轨迹、频谱匹配、CU 功率、卸载波束、卸载时长及计算频率。下面构造新的感知发射波束和接收波束，并证明所得方案仍可行且目标值不增。

首先，任取 $i$ 和 $t$，令
\begin{equation}\label{eq:new-cov}
\mathbf R_i(t)=\sum_{j=1}^{J}\eta_{i,j}(t)p_j(t)
\bm h_{c_j,u_i}^{H}(t)\bm h_{c_j,u_i}(t)+\sigma^2\mathbf I.
\end{equation}
由于 $\sigma^2>0$ 且 $p_j(t)\ge0$，有 $\mathbf R_i(t)\succ0$。将式~\eqref{eq:A} 代入感知 SINR，并利用式~\eqref{eq:cs} 的 Cauchy--Schwarz 不等式，可得对任意单位范数接收波束 $\bm g_i(t)$，
\begin{align}\label{eq:new-sinr-bound}
\operatorname{SINR}_i(t)
&=\xi_0\rho d_i^{-4}(t)\left|\bm a_t^H(\theta_i(t))\tilde{\bm w}_i(t)\right|^2
\frac{\left|\bm g_i^H(t)\bm a_r(\theta_i(t))\right|^2}
{\bm g_i^H(t)\mathbf R_i(t)\bm g_i(t)}\notag\\
&\le\xi_0\rho d_i^{-4}(t)\left|\bm a_t^H(\theta_i(t))\tilde{\bm w}_i(t)\right|^2
\bm a_r^H(\theta_i(t))\mathbf R_i^{-1}(t)\bm a_r(\theta_i(t)).
\end{align}
选择 $\bm g_i^{\rm MVDR}(t)$ 为式~\eqref{eq:mvdr} 的单位范数向量时，式~\eqref{eq:new-sinr-bound} 取等。若发射目标增益为零，各接收波束的感知 SINR 均为零，式~\eqref{eq:mvdr} 仍为一个最大 SINR 的接收波束。

记原方案的感知 SINR 为 $\tilde s_i(t)$，在保持 $\tilde{\bm w}_i(t)$ 不变时，将原接收波束替换为 $\bm g_i^{\rm MVDR}(t)$ 后的感知 SINR 为 $\hat s_i(t)$。由式~\eqref{eq:new-sinr-bound} 可知
\begin{equation}\label{eq:new-ineq}
\hat s_i(t)\ge\tilde s_i(t)\ge0.
\end{equation}
构造新的感知发射波束 $\bm w'_i(t)=c_i(t)\tilde{\bm w}_i(t)$，其中
\begin{equation}\label{eq:new-scale}
c_i(t)=
\begin{cases}
\sqrt{\tilde s_i(t)/\hat s_i(t)},&\hat s_i(t)>0,\\
0,&\hat s_i(t)=0.
\end{cases}
\end{equation}
由式~\eqref{eq:new-ineq}，$0\le c_i(t)\le1$；当 $\hat s_i(t)=0$ 时，必有 $\tilde s_i(t)=0$。故新方案的感知 SINR $s'_i(t)$ 和发射功率满足
\begin{align}\label{eq:new-same-sinr}
s'_i(t)&=c_i^2(t)\hat s_i(t)=\tilde s_i(t),\\
\|\bm w'_i(t)\|^2&=c_i^2(t)\|\tilde{\bm w}_i(t)\|^2\le\|\tilde{\bm w}_i(t)\|^2.\label{eq:new-power}
\end{align}
对所有 UAV 和时隙分别执行上述构造，所得方案记为 $\bm\zeta'$。接下来验证 $\bm\zeta'$ 满足 $P0$ 的约束。

首先，$\bm g_i^{\rm MVDR}(t)$ 满足单位范数约束；式~\eqref{eq:new-power} 表明感知发射功率上限仍成立。根据式~\eqref{eq:new-same-sinr}，感知 SINR 阈值仍满足，且式~\eqref{eq:rate} 给出的感知速率和数据量不变，即
\begin{equation}\label{eq:new-same-rate}
R_i^{\rm sen\prime}(t)=\tilde R_i^{\rm sen}(t),\qquad
\bar D^{\rm sen}R_i^{\rm sen\prime}(t)=\bar D^{\rm sen}\tilde R_i^{\rm sen}(t).
\end{equation}
由于卸载波束、卸载时长和计算频率均保持原值，根据原可行解满足式~\eqref{eq:impact-off} 和~\eqref{eq:impact-delay}，有
\begin{align}
\bar D^{\rm sen}R_i^{\rm sen\prime}(t)
&=\bar D^{\rm sen}\tilde R_i^{\rm sen}(t)
\le \tilde D_{u_i}^{\rm off}(t)\tilde R_i^{\rm off}(t),\label{eq:new-off}\\
\bar D^{\rm sen}+\tilde D_{u_i}^{\rm off}(t)
+\frac{C^{\rm sen}\bar D^{\rm sen}R_i^{\rm sen\prime}(t)}{\tilde f_{u_i}(t)}
&=\bar D^{\rm sen}+\tilde D_{u_i}^{\rm off}(t)
+\frac{C^{\rm sen}\bar D^{\rm sen}\tilde R_i^{\rm sen}(t)}{\tilde f_{u_i}(t)}
\le D_{\max}^{\rm sen}.\label{eq:new-delay}
\end{align}
因此，感知数据卸载约束和计算时延约束仍成立，且 BS 感知计算能耗不变。

其次，对于与 UAV $u_i$ 共享频谱的 CU $c_j$，感知期间 UAV 对 BS 的干扰满足
\begin{equation}\label{eq:new-interference}
\left|\bm h_{u_i,\mathrm{BS}}^H(t)\bm w'_i(t)\right|^2
=c_i^2(t)\left|\bm h_{u_i,\mathrm{BS}}^H(t)\tilde{\bm w}_i(t)\right|^2
\le\left|\bm h_{u_i,\mathrm{BS}}^H(t)\tilde{\bm w}_i(t)\right|^2.
\end{equation}
因此，CU 的第一阶段卸载速率不降低，即 $R_{i,j}^{(1)\prime}(t)\ge\tilde R_{i,j}^{(1)}(t)$。第二、第三阶段速率及各阶段时长保持不变，因此 CU 成功卸载约束仍成立。其余变量和约束均未改变，包括跨时隙轨迹约束，故 $\bm\zeta'$ 是 $P0$ 的可行解。

令 $E_{P0}(\cdot)$ 表示 $P0$ 的长期平均加权能耗。根据式~\eqref{eq:new-same-rate} 和~\eqref{eq:new-power}，BS 计算能耗、CU 能耗、UAV 卸载能耗与飞行能耗均不变，而感知发射能耗不增加。因此
\begin{align}\label{eq:new-objective}
E_{P0}(\bm\zeta')-E_{P0}(\tilde{\bm\zeta})
&=-\frac{\omega_2\bar D^{\rm sen}}{t_{\max}}
\sum_{t=1}^{t_{\max}}\sum_{u_i\in\mathcal U}
\bigl(1-c_i^2(t)\bigr)\left\|\tilde{\bm w}_i(t)\right\|^2
\le0.
\end{align}

最后，令 $\bm\zeta^\star$ 为 $P0$ 的一个全局最优解，$E_{P0}^\star=E_{P0}(\bm\zeta^\star)$ 为其最优值。对 $\bm\zeta^\star$ 执行上述构造，得到使用 MVDR 接收波束的可行解 $\bm\zeta^{\rm MVDR}$。由式~\eqref{eq:new-objective}，有
\begin{equation}\label{eq:new-opt-upper}
E_{P0}(\bm\zeta^{\rm MVDR})\le E_{P0}(\bm\zeta^\star)=E_{P0}^\star.
\end{equation}
另一方面，由于 $\bm\zeta^{\rm MVDR}$ 对 $P0$ 可行，且 $E_{P0}^\star$ 为全局最小值，有
\begin{equation}\label{eq:new-opt-lower}
E_{P0}(\bm\zeta^{\rm MVDR})\ge E_{P0}^\star.
\end{equation}
结合式~\eqref{eq:new-opt-upper} 和~\eqref{eq:new-opt-lower}，得到
\begin{equation}
E_{P0}(\bm\zeta^{\rm MVDR})=E_{P0}^\star.
\end{equation}
因此，$\bm\zeta^{\rm MVDR}$ 是 $P0$ 的一个全局最优解，定理得证。\hfill$\square$

